\section*{Chapitre \ref{ch:m1:volatility-first-contact} --- La volatilité comme grandeur négociée~: premier contact}

\begin{solution}{exo:m1:volatility-first-contact:1}
$F = 100 + \mathrm{e}^{0.01}(4.06 - 3.87) = 100.19$. Carry implicite~: $100 -
\mathrm{e}^{-0.01} \times 100.19 = 0.81$~: un dividende attendu d'environ quatre-vingts
centimes (un strike, avec des prix arrondis au centime, est précis à un ou deux centimes près).
\end{solution}

\begin{solution}{exo:m1:volatility-first-contact:2}
Le strike est légèrement inférieur au forward, donc le call est un peu dans la
monnaie et le put un peu hors de la monnaie~: on prend leur moyenne, 3.965. $\sigma \approx
3.965/(0.4 \times 100.2 \times 0.990 \times 0.5) = 20.0\%$, contre 20.1 obtenu avec
le solveur.
\end{solution}

\begin{solution}{exo:m1:volatility-first-contact:3}
$2.55 \times \$1\,000 \times 40 = \$102\,000$~; $2.55/0.05 = 51$ ticks.
\end{solution}

\begin{solution}{exo:m1:volatility-first-contact:4}
La borne inférieure du call est $\mathrm{e}^{-rT}(F - K) = 0.990 \times 10.20 =
10.10$. Un prix de 10.05 est inférieur à cette borne~: aucune volatilité, aussi faible soit-elle, ne le produit.
Le système doit signaler la cotation (périmée, croisée avec le cours de l'action, ou
le forward est erroné, par exemple un dividende pas encore pris en compte dans les données) et l'exclure
de toute surface~; renvoyer zéro ou la valeur précédente masque le problème.
\end{solution}

\begin{solution}{exo:m1:volatility-first-contact:5}
Carry $= 100 - \mathrm{e}^{-0.01} \times 98.10 = 2.88$ sur un trimestre~: 11.5~\% par
an. Sans dividende, c'est le \omterm{def:m1:financing:seclending}{frais d'emprunt} que paie l'arbitragiste marginal pour
être vendeur à découvert~: l'action est un special, et les puts paraissent chers par rapport
aux calls, précisément de ce montant. Il ne s'agit pas d'un arbitrage pour quiconque doit payer
le même frais.
\end{solution}

\begin{solution}{exo:m1:volatility-first-contact:6}
$\sum r^2 = 0.16 + 1.21 + 0.49 + 4.00 + 2.25 = 8.11$ (en \%$^2$)~; moyenne 1.622~;
volatilité quotidienne 1.274~\%~; annualisée $\times\sqrt{252}$~: 20.2~\%. Sur ces
cinq jours, l'acheteur d'options à 16 avait raison~; cinq jours sont bien trop peu
pour en dire davantage.
\end{solution}

\begin{solution}{exo:m1:volatility-first-contact:7}
Avec des prix au centime~: forward moyen de 100.20, écart strike par strike inférieur à trois
centimes, dividende implicite de 0.80, skew de 6.2 points. Avec des prix à cinq centimes, le
forward et le dividende bougent de quelques centimes (l'estimation du dividende, 0.80, a une
erreur de même taille que l'arrondi, soit plusieurs pour cent de sa propre valeur), tandis que les
volatilités proches de la monnaie changent à peine~; les volatilités des options les moins
chères, valant 0.05 à 0.35, deviennent inutilisables, puisque cinq centimes représentent une
part importante de leur prix. Les marchés larges brouillent d'abord les ailes.
\end{solution}

\begin{solution}{exo:m1:volatility-first-contact:8}
L'indice ne peut pas être acheté~; le contrat à terme, lui, peut l'être, à 15.6, et la courbe
contient déjà le retour à la moyenne~: le deuxième contrat est à 16.7, le septième à
19.0. Si l'indice reste à 14, le contrat acheté à 15.6 converge vers 14 et
perd 1.6 point, soit 1\,600~\$ par contrat, en un mois, chaque mois. Acheter revient à
parier que l'indice monte \emph{plus vite que ne l'indique la courbe}, face à un carry que
les vendeurs de volatilité encaissent le reste du temps.
\end{solution}

\begin{solution}{pb:m1:volatility-first-contact:1}
\textbf{1.} $F_K = K + \mathrm{e}^{0.01}(C - P)$~: 100.20, 100.20, 100.19,
100.20, 100.20.
\textbf{2.} Les prix sont arrondis au centime puis multipliés par 1.01. Les strikes les moins
fiables sont ceux où l'une des deux options est loin dans la monnaie~:
sa cotation est large et elle peut comporter une prime d'exercice anticipé si elle est américaine.
\textbf{3.} 100.20.
\textbf{4.} $100 - 0.990 \times 100.20 = 0.80$.
\textbf{5.} Non~: le carry implicite est égal au dividende attendu (0.80, soit
0.797 aujourd'hui). Il ne reste rien pour un \omterm{def:m1:financing:seclending}{frais d'emprunt}.
\textbf{6.} L'option hors de la monnaie~: le put sous le forward, le call
au-dessus. Sa cotation est la plus serrée par rapport à sa sensibilité à la volatilité,
elle n'a pas de \omterm{def:m1:option-contracts-and-exchanges:series}{valeur intrinsèque} à retrancher, et ne comporte guère de prime d'exercice anticipé.
\textbf{7.} $23.8 - 17.6 = 6.2$ points~; 3.1 points pour 10~\% du strike.
\textbf{8.} $(4.06 + 3.87)/2 = 3.965$~; $3.965/(0.4 \times 100.2 \times 0.990
\times 0.5) = 20.0\%$.
\textbf{9.} Avec $F = 100$ à la place de 100.20, la \omterm{def:m1:option-contracts-and-exchanges:series}{valeur intrinsèque} du call est
sous-estimée de 0.20, donc davantage de ses 11.22 est attribué à la volatilité~: 25.1
au lieu de 23.8. Les options dans la monnaie ont peu de sensibilité à la volatilité, si bien
qu'une petite erreur sur le forward devient une grande erreur sur la \omterm{def:m1:volatility-first-contact:iv}{volatilité implicite}~;
et le call et le put au même strike ne concordent alors plus, ce qui
est le symptôme à rechercher.
\textbf{10.} $4.03 - 3.90 - 100.01 + 0.797 + 99.005 = -0.08$.
\textbf{11.} $3.84 - 4.09 + 99.99 - 0.797 - 0.05 - 99.005 = -0.11$.
\textbf{12.} Les deux perdent~: la \omterm{prop:m1:volatility-first-contact:parity}{parité call-put} tient à l'intérieur des fourchettes, et la
chaîne est cohérente avec un seul forward. C'est l'état normal, entretenu par
les desks pour lesquels ces deux nombres sont les plus proches de zéro.
\textbf{13.} Une conversion donnerait $4.60 - 3.90 - 100.01 + 0.797 + 99.005 =
+0.49$ par action. Il faut d'abord vérifier que la cotation est réelle et à jour (taille,
bourse, heure), que le cours de l'action est à jour, qu'aucun dividende ni
aucune opération sur titre n'a modifié le contrat, et que l'option ne va pas être
ajustée~; on la vend ensuite.
\textbf{14.} Pas par le prix. Le dividende restant inchangé, le forward vaut
97.17, et 2.74 correspond à une \omterm{def:m1:volatility-first-contact:iv}{volatilité implicite} de 20.5 contre 20.1 la veille~:
en termes de volatilité, le call est légèrement plus cher. On compare les volatilités,
et au-delà, les volatilités à moneyness comparable.
\textbf{15.} Parce que la volatilité est stable lorsque l'action se déplace, comparable
d'un strike à l'autre, d'une échéance ou d'un sous-jacent à l'autre, et qu'elle est la grandeur
à laquelle une position d'options couverte est exposée.
\textbf{16.} Que le marché paie davantage, par unité de risque de modèle, pour
se protéger contre les baisses que pour participer aux hausses~: crainte de sauts
à la baisse, demande de puts de la part des détenteurs de l'action, offre de calls
de la part des vendeurs couverts.
\textbf{17.} Il pondère les options par $1/K^2$ et englobe toute la bande~: les
puts coûteux aux bas strikes comptent le plus, si bien que le nombre est une moyenne
du smile, penchée vers son côté haut.
\textbf{18.} Les vendeurs d'options supportent un risque qui se paie mal dans les mauvaises périodes~;
ils sont rémunérés par une prime, comme un assureur. Une \omterm{def:m1:volatility-first-contact:iv}{volatilité implicite} supérieure à la réalisée en
moyenne correspond à cette prime, et elle est restituée, avec intérêts, dans des mois comme
février 2018.
\textbf{19.} \textbf{Un dividende implicite de 0.80 et un skew 90--110 de 6.2
points de volatilité.}
\textbf{20.} Le prix d'une option, exprimé dans l'unité qui reste
fixe lorsque le sous-jacent bouge.
\end{solution}

\begin{solution}{iq:m1:volatility-first-contact:1}
$C - P = \mathrm{e}^{-rT}(F - K)$ pour des options européennes de même strike et de même
échéance. Elle suppose seulement que l'on puisse négocier le forward~: acheter ou vendre à
découvert le sous-jacent, et emprunter ou prêter des liquidités. Les dividendes et les \omterm{def:m1:financing:seclending}{frais
d'emprunt} entrent par $F$. Aucun modèle, aucune distribution. Pour les options américaines, elle
devient un couple d'inégalités.
\iqlookfor{indépendante de tout modèle~; le forward, non le spot~; la réserve sur les options américaines.}
\end{solution}

\begin{solution}{iq:m1:volatility-first-contact:2}
La volatilité qui fait coïncider la formule de Black--Scholes avec le prix de marché~:
une conversion d'unités. Elle contient la volatilité attendue, une prime de risque
pour la supporter, ainsi que l'offre et la demande pour ce strike~; elle varie avec le strike,
ce qu'aucune prévision unique ne saurait faire. En moyenne, elle a dépassé la \omterm{def:m1:volatility-first-contact:realised}{volatilité
réalisée} ultérieure.
\iqlookfor{la prime de risque, et le smile comme preuve qu'il ne s'agit pas d'une
prévision.}
\end{solution}

\begin{solution}{iq:m1:volatility-first-contact:3}
Vérifier la prime par rapport à ses bornes (\omterm{def:m1:option-contracts-and-exchanges:series}{valeur intrinsèque} actualisée et
forward ou strike actualisé), et la rejeter si elle les sort. Utiliser l'option hors de
la monnaie via la parité. La bissection converge toujours~; Newton est plus rapide, mais
échoue là où le vega est très faible, loin dans ou hors de la monnaie~; il faut donc le
protéger par des encadrements. Attention au mauvais forward (dividendes), à la prime
américaine, aux options le jour de l'échéance avec $T \to 0$, aux prix d'un tick lorsque
la réponse couvre vingt points de volatilité, et aux conventions de temps (temps
calendaire ou temps ouvré).
\iqlookfor{contrôle des bornes, effondrement du vega, et le forward comme première source
d'erreur.}
\end{solution}

\begin{solution}{iq:m1:volatility-first-contact:4}
Un dividende que je n'ai pas inclus, ou plus grand que celui que j'avais supposé. Des \omterm{def:m1:financing:seclending}{frais
d'emprunt}~: les actions difficiles à emprunter ont des forwards bas. Le taux d'intérêt que
j'utilise. L'\omterm{def:m1:option-contracts-and-exchanges:style}{exercice américain}~: le put porte une prime d'exercice anticipé. Des cotations
périmées ou larges, comparées au milieu. Une opération sur titres annoncée qui modifie
le sous-jacent livrable. Ce n'est un arbitrage que si la conversion ou la réversion
est rentable aux prix auxquels je peux réellement négocier, emprunt compris.
\iqlookfor{dividendes et emprunt avant toute chose.}
\end{solution}

\begin{solution}{iq:m1:volatility-first-contact:5}
À partir des options d'indice hors de la monnaie de deux échéances encadrant
trente jours~: la \omterm{def:m1:volatility-first-contact:realised}{variance} de chaque échéance est une somme de prix milieu pondérés
par $\Delta K/K^2$, avec une petite correction pour la distance entre le forward et
le strike le plus proche~; les deux \omterm{def:m1:volatility-first-contact:realised}{variances} sont interpolées à trente jours, et
l'indice vaut cent fois la racine carrée. C'est le prix d'un portefeuille qui devrait être
rééquilibré en continu et dont la composition change à chaque instant~; on peut négocier
des contrats à terme et des options sur lui, qui évaluent sa valeur future, et non
l'indice lui-même.
\iqlookfor{le panier et ses pondérations~; les contrats à terme ne sont pas l'indice.}
\end{solution}

\begin{solution}{iq:m1:volatility-first-contact:6}
Des produits vendeurs de volatilité à court terme, qui avaient encaissé le roll pendant
des années, ont pris beaucoup d'ampleur. Ils se rééquilibrent chaque jour pour maintenir
un effet de levier constant~: après une hausse des contrats à terme, les produits inverses
comme les produits à levier acheteur doivent acheter, dans les dernières minutes,
un montant proportionnel au mouvement. Le 5 février, une baisse de 4~\% de l'indice
a fait monter les contrats à terme~; les achats requis étaient importants par rapport
à la liquidité de ces contrats, ont poussé ceux-ci encore plus haut, et un produit inverse
a perdu presque tout en une séance. Visibles à l'avance~: la taille des produits par rapport
au volume des contrats à terme, la convexité de la règle de rééquilibrage (la demande
croît avec le mouvement), la concentration de cette demande à un moment de la journée, et une position dont le pire jour dépasse son profit historique cumulé.
\iqlookfor{l'arithmétique du rééquilibrage et la taille relative à la liquidité, et non
simplement «~la volatilité a grimpé~».}
\end{solution}
